\documentclass[../DoAn.tex]{subfiles}
\begin{document}

Chương này trình bày các nền tảng kỹ thuật và công nghệ chính được sử dụng để xây dựng hệ thống đặt vé và quản lý sự kiện trên thiết bị di động. Do đặc thù của bài toán yêu cầu đồng thời hỗ trợ nhiều nghiệp vụ như quản lý tài khoản, tổ chức sự kiện, đặt vé, thông báo, nhắc lịch, trò chuyện và kiểm soát check-in, giải pháp kỹ thuật cần bảo đảm tính linh hoạt khi phát triển, khả năng mở rộng hợp lý và thuận tiện cho việc tích hợp nhiều chức năng trên cùng một hệ thống. Vì vậy, đồ án lựa chọn kiến trúc client-server theo hướng mobile-first, trong đó ứng dụng di động đảm nhiệm tương tác với người dùng, còn máy chủ backend chịu trách nhiệm xử lý nghiệp vụ và lưu trữ dữ liệu tập trung.

\section{Kiến trúc ứng dụng di động kết hợp dịch vụ backend}
Kiến trúc tổng thể của hệ thống được xây dựng theo mô hình ứng dụng khách và máy chủ. Mô hình này được lựa chọn để giải quyết yêu cầu đồng bộ dữ liệu giữa nhiều nhóm người dùng như người tham dự, người tổ chức và quản trị viên, đồng thời hỗ trợ quản lý nghiệp vụ tập trung như xác thực, quản lý sự kiện, quản lý vé, xử lý thông báo và kiểm soát quyền truy cập. Với bài toán của đề tài, một số hướng tiếp cận thay thế có thể kể đến như ứng dụng web thuần, ứng dụng desktop hoặc mô hình backendless sử dụng hoàn toàn dịch vụ bên thứ ba. Tuy nhiên, ứng dụng web thuần thường kém phù hợp hơn trong trải nghiệm di động chuyên biệt, ứng dụng desktop không phù hợp với bối cảnh sử dụng thường xuyên khi tham gia sự kiện, còn mô hình backendless tuy rút ngắn thời gian phát triển nhưng làm giảm khả năng kiểm soát nghiệp vụ tùy biến. Vì vậy, mô hình ứng dụng di động kết hợp backend riêng được xem là phù hợp hơn cho yêu cầu của hệ thống.

\section{React Native và Expo cho lớp ứng dụng khách}
Phía ứng dụng khách của hệ thống được phát triển bằng React Native kết hợp với Expo. React Native cho phép xây dựng ứng dụng di động đa nền tảng bằng JavaScript/TypeScript, đồng thời cung cấp khả năng tái sử dụng lớn về cấu trúc giao diện và logic hiển thị. Expo hỗ trợ đáng kể trong quá trình phát triển thông qua hệ sinh thái thư viện sẵn có cho thông báo, định vị, camera, quét mã và nhiều tính năng thường gặp trong ứng dụng di động hiện đại. Hai công nghệ này được sử dụng để giải quyết yêu cầu xây dựng nhanh một ứng dụng mobile-first có nhiều màn hình nghiệp vụ như đăng nhập, khám phá sự kiện, đặt vé, quản lý hồ sơ, thông báo, nhắn tin và check-in.

Các lựa chọn thay thế cho lớp ứng dụng khách có thể là native Android/iOS, Flutter hoặc Progressive Web App. Phát triển native mang lại khả năng tối ưu cao nhưng làm tăng chi phí phát triển song song cho hai nền tảng. Flutter có lợi thế về hiệu năng giao diện, tuy nhiên trong phạm vi đề tài, React Native kết hợp Expo phù hợp hơn nhờ hệ sinh thái gần gũi với JavaScript, dễ tích hợp với các thư viện hiện có và rút ngắn thời gian hiện thực hóa sản phẩm thử nghiệm.

\section{Node.js và Express cho lớp dịch vụ nghiệp vụ}
Phía máy chủ của hệ thống sử dụng Node.js và Express. Node.js phù hợp với các ứng dụng cần xử lý nhiều kết nối đồng thời và có số lượng lớn yêu cầu I/O như truy vấn dữ liệu, xác thực người dùng, gửi thông báo và trao đổi thời gian thực. Express được lựa chọn để xây dựng các \gls{API} \gls{REST} cho những nghiệp vụ chính như quản lý tài khoản, sự kiện, vé, thanh toán, nhắc lịch, bookmark, đánh giá và trò chuyện, với dữ liệu trao đổi chủ yếu ở định dạng \gls{JSON}. Cặp công nghệ này giải quyết yêu cầu xây dựng một backend thống nhất, dễ tổ chức theo mô hình route-controller-service và thuận lợi trong việc mở rộng chức năng sau này.

Một số phương án thay thế có thể dùng cho backend là Spring Boot, Django hoặc NestJS. Spring Boot mạnh về cấu trúc enterprise nhưng tương đối nặng cho phạm vi đồ án ứng dụng cỡ vừa. Django cung cấp nhiều thành phần tích hợp sẵn nhưng không đồng bộ về ngôn ngữ với phía ứng dụng khách. NestJS là lựa chọn hiện đại hơn trong hệ sinh thái Node.js, tuy nhiên Express được ưu tiên trong đề tài vì đơn giản, phổ biến và phù hợp với mức độ kiểm soát trực tiếp mà nhóm phát triển mong muốn.

\section{MongoDB và Mongoose cho lưu trữ dữ liệu}
Hệ thống sử dụng MongoDB làm hệ quản trị cơ sở dữ liệu và Mongoose làm thư viện ánh xạ dữ liệu cho Node.js. Việc lựa chọn này nhằm giải quyết yêu cầu lưu trữ dữ liệu có cấu trúc linh hoạt như người dùng, sự kiện, hạng vé, vé điện tử, thanh toán, phòng chat, tin nhắn và thông báo. Trong bài toán sự kiện, nhiều thực thể có thể thay đổi về thuộc tính theo từng phiên bản phát triển, chẳng hạn sự kiện có thể mở rộng thêm hạng vé, trạng thái duyệt hay thông tin tọa độ. Mô hình tài liệu của MongoDB, vốn thuộc nhóm \gls{NoSQL}, giúp thích nghi tốt với sự thay đổi đó, trong khi Mongoose hỗ trợ định nghĩa schema, kiểm tra dữ liệu, quan hệ tham chiếu và thao tác với cơ sở dữ liệu thuận tiện hơn.

Các lựa chọn thay thế có thể kể đến như MySQL, PostgreSQL hoặc Firebase Firestore. Hệ quản trị quan hệ phù hợp với dữ liệu chuẩn hóa chặt chẽ nhưng có thể làm tăng độ phức tạp khi mô hình nghiệp vụ cần mở rộng linh hoạt trong giai đoạn thử nghiệm. Firestore hỗ trợ thời gian thực tốt nhưng việc tự kiểm soát toàn bộ nghiệp vụ phía máy chủ sẽ hạn chế hơn. Vì vậy, MongoDB kết hợp Mongoose là lựa chọn phù hợp giữa tính linh hoạt phát triển và khả năng kiểm soát dữ liệu.

\section{JWT cho xác thực và phân quyền}
Để giải quyết yêu cầu xác thực người dùng và bảo vệ các tài nguyên nghiệp vụ của hệ thống, đề tài sử dụng cơ chế xác thực dựa trên \gls{JWT} \cite{jwt2026rfc}. Sau khi đăng nhập thành công, máy chủ cấp token cho người dùng, và token này được gửi kèm trong các yêu cầu tới những \gls{API} cần xác thực. Cách tiếp cận này phù hợp với ứng dụng di động vì cho phép duy trì trạng thái đăng nhập tương đối gọn nhẹ, không phụ thuộc phiên làm việc lưu trên máy chủ theo kiểu session truyền thống. Ngoài xác thực, hệ thống còn áp dụng phân quyền theo vai trò người dùng và quản trị viên để kiểm soát các chức năng quản trị hoặc nghiệp vụ chuyên biệt.

Các phương án thay thế cho JWT là session-based authentication hoặc OAuth hoàn chỉnh từ bên thứ ba. Session phù hợp với web truyền thống nhưng kém linh hoạt hơn trong môi trường ứng dụng di động phân tán. OAuth rất mạnh với kịch bản đăng nhập liên thông, tuy nhiên vượt quá phạm vi cốt lõi của đề tài. Do đó, JWT được lựa chọn vì đáp ứng tốt yêu cầu xác thực API, đơn giản trong triển khai và phù hợp với kiến trúc hệ thống hiện tại.

\section{Socket.IO cho giao tiếp thời gian thực}
Đối với yêu cầu trao đổi tin nhắn và cập nhật trạng thái phòng chat gần thời gian thực, hệ thống sử dụng Socket.IO. Công nghệ này cho phép thiết lập kết nối hai chiều giữa máy khách và máy chủ, phù hợp với bài toán gửi tin nhắn, tham gia phòng trò chuyện và phát sự kiện cập nhật ngay khi có nội dung mới. Trong bối cảnh ứng dụng sự kiện, tính năng giao tiếp trực tiếp giữa người dùng với người tổ chức hoặc giữa các thành viên liên quan làm tăng khả năng tương tác và hỗ trợ xử lý nhanh các tình huống phát sinh.

Các lựa chọn thay thế có thể là polling định kỳ, Server-Sent Events hoặc dịch vụ chat của bên thứ ba. Polling dễ triển khai nhưng gây lãng phí tài nguyên và độ trễ cao hơn. Server-Sent Events chủ yếu phù hợp với luồng dữ liệu một chiều từ máy chủ xuống máy khách. Dịch vụ bên thứ ba rút ngắn thời gian triển khai nhưng làm giảm khả năng tùy biến nghiệp vụ. Vì vậy, Socket.IO là lựa chọn phù hợp hơn cho chức năng chat của hệ thống.

\section{Thông báo đẩy và nhắc lịch sự kiện}
Hệ thống tích hợp Expo Notifications để hỗ trợ thông báo đẩy trên thiết bị di động. Công nghệ này được sử dụng để giải quyết yêu cầu gửi thông báo khi có tin nhắn mới, xác nhận thanh toán, thay đổi liên quan tới sự kiện và nhắc lịch trước thời điểm diễn ra sự kiện. Bên cạnh đó, backend triển khai bộ lập lịch kiểm tra các sự kiện sắp diễn ra và gửi thông báo trước các mốc thời gian nhất định. Đây là thành phần quan trọng đối với bài toán sự kiện, bởi người dùng không chỉ cần mua vé mà còn cần được nhắc đúng lúc để tránh bỏ lỡ chương trình đã đăng ký.

Một số lựa chọn thay thế có thể là Firebase Cloud Messaging kết hợp giải pháp native riêng hoặc các nền tảng thông báo thương mại. Tuy nhiên, Expo Notifications thuận lợi hơn trong phạm vi đề tài vì tích hợp tốt với hệ sinh thái Expo, giảm đáng kể khối lượng cấu hình và phù hợp với quy mô hệ thống thử nghiệm.

\section{Mã QR trong quản lý vé điện tử}
Để giải quyết yêu cầu số hóa vé và kiểm soát check-in tại sự kiện, đề tài sử dụng cơ chế sinh và đọc mã \gls{QR} cho từng vé điện tử thông qua thư viện tạo mã phù hợp trên nền Node.js. Mỗi vé sau khi thanh toán thành công được gắn một chuỗi dữ liệu định danh, từ đó hỗ trợ kiểm tra tính hợp lệ khi người tổ chức thực hiện check-in. Hướng tiếp cận này giúp giảm phụ thuộc vào vé giấy, rút ngắn thời gian xác nhận tham dự và tạo trải nghiệm thuận tiện hơn cho cả người dùng lẫn người tổ chức.

Những phương án thay thế cho kiểm soát check-in có thể là mã vạch một chiều, danh sách thủ công hoặc xác nhận bằng thông tin định danh cá nhân. So với các phương án này, mã QR có ưu điểm là dễ sinh tự động, dễ quét trên thiết bị di động và phù hợp với ngữ cảnh triển khai nhanh trong môi trường ứng dụng sự kiện.

\section{Hỗ trợ sinh nội dung sự kiện bằng AI}
Ngoài các nghiệp vụ cốt lõi, hệ thống còn tích hợp khả năng gợi ý nội dung sự kiện bằng mô hình ngôn ngữ lớn thông qua dịch vụ \gls{API}. Thành phần này được sử dụng để hỗ trợ người tổ chức khi cần đề xuất tiêu đề, mô tả ngắn và các hạng vé ban đầu cho sự kiện. Trong bài toán thực tế, nhiều người tổ chức gặp khó khăn khi phải đồng thời chuẩn bị nội dung quảng bá và cấu hình chương trình. Việc bổ sung chức năng hỗ trợ bằng \gls{AI} không thay thế vai trò của người dùng, nhưng giúp rút ngắn thời gian chuẩn bị và tăng tính thuận tiện cho thao tác tạo sự kiện.

Các hướng thay thế cho bài toán này có thể là nhập tay hoàn toàn, sử dụng mẫu nội dung cố định hoặc tích hợp một công cụ gợi ý rule-based. Tuy nhiên, nhập tay hoàn toàn làm tăng thời gian thao tác, trong khi mẫu cố định thường thiếu linh hoạt. Vì vậy, hướng tích hợp AI được lựa chọn như một thành phần hỗ trợ mở rộng, phù hợp với xu thế cá nhân hóa nội dung trong các ứng dụng hiện đại.

\section{Kết chương}
Chương này đã trình bày các nền tảng kỹ thuật và công nghệ chính được sử dụng trong hệ thống đặt vé và quản lý sự kiện của đồ án. Mỗi lựa chọn công nghệ đều được xem xét dựa trên yêu cầu nghiệp vụ tương ứng, đồng thời đối chiếu với một số phương án thay thế để làm rõ lý do lựa chọn. Từ việc sử dụng kiến trúc client-server theo hướng mobile-first, React Native và Expo cho ứng dụng khách, Node.js và Express cho dịch vụ backend, MongoDB cho lưu trữ dữ liệu, JWT cho xác thực, Socket.IO cho thời gian thực, thông báo đẩy cho nhắc lịch, mã QR cho vé điện tử và AI cho hỗ trợ tạo nội dung, hệ thống đã hình thành được nền tảng kỹ thuật tương đối đầy đủ để triển khai các chức năng cốt lõi của bài toán. Trên cơ sở này, chương tiếp theo có thể tập trung trình bày chi tiết hơn về thiết kế, hiện thực và đánh giá hệ thống.


\end{document}
